\chapter{Contexto}
\label{cap:contexto}

En el contexto se presenta el ámbito en el que se ha desarrollado el TFM.
Considerar hablar desde aspectos más generales, a aquellos más
particulares relacionados con el TFM. Finalizar con un párrafo que hable de la problemática en la cual se ha focalizado el TFM detallando la necesidad de profundizar en esa área de investigación.

\section{Marco Macroscópico: La Descentralización del Sector Energético}
\label{sec:marco_macroscopico}

El presente Trabajo de Fin de Máster se contextualiza en el escenario global de la transición energética y la digitalización industrial, un proceso caracterizado por la migración activa desde un modelo de generación eléctrica centralizado basado en combustibles fósiles hacia un ecosistema distribuido y bajo en carbono. Este cambio de paradigma está impulsado por políticas internacionales de descarbonización y por planes estratégicos que imponen un despliegue masivo de fuentes de energía renovables \cite{aee2026}.Las directrices globales de descarbonización se rigen por el marco jurídico internacional del Acuerdo de París \cite{unfcccParisAgreementPub}, el cual establece el límite del incremento térmico en 1.5~°C. Los escenarios de descarbonización global exigen triplicar la capacidad instalada de energías renovables y optimizar las redes de distribución para absorber la generación distribuida, fijando marcas hacia el año 2030 \cite{ieaNetZero2023Update}.

En España el proceso de transición energética se enmarca dentro del Plan Nacional Integrado de Energía y Clima (PNIEC) 2021--2030 \cite{mitecoPNIEC2023}, el cual establece objetivos vinculantes para la reducción de emisiones de gases de efecto invernadero (GEI) en un 23\% para el año 2030 respecto a los niveles de 1990, con una meta específica de alcanzar una potencia instalada de energías renovables de 74~GW. Requiriendo una gestión eficiente que garantice la estabilidad y seguridad del sistema eléctrico nacional.\par

\section{Marco Microscópico: La Complejidad de la Gestión de Parques Eólicos}
Desde la perspectiva de la automatización, la transición hacia las energías renovables introduce una complejidad de campo. Mientras que una central térmica tradicional concentra una elevada potencia en un único punto monolítico, un parque eólico moderno fragmenta dicha potencia en decenas de activos individuales (aerogeneradores) distribuidos a lo largo de amplias extensiones geográficas. Cada una de estas turbinas constituye un sistema ciberfísico (CPS) complejo, dotado de su propia instrumentación, electrónica de potencia y lazos de control internos en tiempo real (regulación de orientación, paso de pala y par del generador). En adición la integración de activos de diferentes fabricantes o generaciones tecnológicas introduce una heterogeneidad de protocolos y estructuras de datos que dificulta la supervisión unificada. Por ello se han desarollado estandares internacionales de interoperabilidad como la norma IEC 61400-25 o la norma IEC 61850 \cite{vattenfall61400} \cite{iec61400-25}.



\section{Ámbito Operativo: La Sala de Control Principal (MCR)}
\label{sec:ambito_operativo_mcr}

Ante la imposibilidad de operar eficientemente cientos de infraestructuras dispersas mediante intervenciones locales, el sector eólico ha centralizado su gestión en las denominadas de Organización y Mantenimiento  (O and M). El ámbito en el que se encuadra este TFM emula de forma directa el entorno de ingeniería de una MCR corporativa, cuyo propósito es unificar bajo una misma capa de software la adquisición de datos y el diagnóstico de múltiples parques eólicos y sus respectivas subestaciones eléctricas de evacuación.
